\documentclass[10pt,a4paper]{amsart}
\usepackage[margin=19mm]{geometry}
\usepackage{amsmath,amssymb,mathtools}
\usepackage[scr=rsfs,cal=euler]{mathalfa}
\usepackage{xfrac}
\usepackage[unicode,hidelinks,bookmarks=false]{hyperref}
\newcommand{\ie}{i.e.\ }
\newcommand{\cf}{cf.\ }
\newcommand{\resp}{respectively\ }
\newcommand{\Subsection}{\subsection}
\providecommand{\nobreakdash}{\nobreak\hbox{-}}
\setlength{\emergencystretch}{2em}
\multlinegap0pt
\setlength{\voffset}{-.25in}
\usepackage{latexsym}
\usepackage{anyfontsize}
\usepackage{tikz}
\usetikzlibrary{cd} %showing a commutative diagram for demonstration purposes; the author should load whatever they need.
\usepackage{color}
\usepackage{url}
\newcommand{\bburl}[1]{\textcolor{blue}{\url{#1}}}
\newcommand{\monthyear}[1]{%
  \def\@monthyear{\uppercase{#1}}}
\newcommand{\volnumber}[1]{%
  \def\@volnumber{\uppercase{#1}}}
\numberwithin{equation}{section} %change this to make globally numbered equations
\newtheorem{thm}{Theorem}[section] %remove [section] to make globally numbered environments
\newtheorem{theorem}[thm]{Theorem}
\newtheorem{lemma}[thm]{Lemma}
\newtheorem{example}[thm]{Example}
\newtheorem{definition}[thm]{Definition}
\newtheorem{proposition}[thm]{Proposition}
\newtheorem{corollary}[thm]{Corollary}
\numberwithin{table}{section} %change this and the following line to make globally numbered tables and figures
\numberwithin{figure}{section}
\newcommand{\helv}{%
  \fontfamily{phv}\fontseries{m}\fontsize{9}{11}\selectfont}
\def\NN{\mathbb{N}}
\def\PP{\mathbb{P}}
\def\ZZ{\mathbb{Z}}
\def\FF{\mathbb{F}}
\def\TT{\mathbb{T}}
\def\QQ{\mathbb{Q}}
\def\RR{\mathbb{R}}
\def\gooddoldprimes{\mathbb{G}_{\rm{D}}}
\def\baddoldprimes{\mathbb{B}_{\rm{D}}}
\def\goodrealizableprimes{\mathbb{G}_{\rm{R}}}
\def\badrealizableprimes{\mathbb{B}_{\rm{R}}}
\newcommand\ord{{\rm{ord}}}
\newcommand\failure{{\rm{Fail}}}
\newcommand\fix{{\rm{Fix}}}
\newcommand\least{{\mathcal{L}}}
\newcommand\fixset{{\rm{FixSet}}}
\newcommand\trace{{\rm{trace}}}
\newcommand\notdivides{\mathrel{\kern-3pt\not\!\kern4.3pt\bigm|}}
\newcommand\divides\mid
\newcommand\smalldivides{\mathrel{\kern-2pt\kern3.5pt|}}
\newcommand\smallnotdivides{\mathrel{\kern-2pt\not\!\kern3.5pt|}}
\def\ge{\geqslant}
\def\le{\leqslant}
\def\eul{\rm{e}}
\DeclareMathOperator{\denom}{denom}
\DeclareMathOperator{\support}{supp}
\begin{document}
\title[The local Dold congruence for Fibonacci and Lucas sequences]{The local Dold congruence for Fibonacci and Lucas sequences}
\author{Sompong Chuysurichay, Sawian Jaidee, Chatchawan Panraksa, Thomas Ward}
\date{}
\maketitle
\noindent\emph{Source and licence.} The local Dold congruence for Fibonacci and Lucas sequences. Version arXiv:2609.24544v1 (CC BY 4.0). This material is distributed under the Creative Commons Attribution 4.0 International licence, http://creativecommons.org/licenses/by/4.0/, as recorded in the arXiv OAI licence metadata for this exact version and shown on the abstract page https://arxiv.org/abs/2609.24544v1. Full rights notice: licenses/arxiv-2609.24544-CC-BY-4.0.txt.\par\smallskip
\noindent\textit{Editorial scope.} Counted unit: the original \texttt{theorem} environment \texttt{theoremLucas} at source lines 516--528 together with its complete proof at lines 530--639; the delivered document keeps source lines 191--640 so that every referenced label and every citation resolves inside it. Native editable source is the paper's own concatenated TeX; the AMS wrapper is editorial formatting only. No publisher class, upstream script or unrelated artwork is redistributed.
\begin{equation}\label{equationDoldForIntegerSequence2}
a_{mq^s}
\equiv
a_{mq^{s-1}}
\pmod{q^s}
\end{equation}
for all~$q\in\PP$,~$m\in\NN\setminus q\mathbb{N}$, and~$s\ge1$.
We refer to the survey by Byszewski {\it{et al.}}~\cite{MR4332826}
for an explanation of this terminology
and to Minton~\cite{MR3195758} for a classification of
these properties for linear recurrence sequences.

An integer sequence may fail to be realizable for
many different reasons, but Miska and Ward~\cite{MR4361581}
explored
a particularly simple notion of \emph{almost realizability},
where the failure can be repaired easily,
which in our context also applies to the Dold property.

\begin{definition}[Almost Dold sequences]\label{definitionAlmostDold}
For an integer sequence~$a$,
write~$\failure(a)$
for the least common multiple of
the denominators of the sequence~$(\frac{1}{n}\least_n(a))_{n\in\NN}$
if this is finite, in which case~$a$
is called \emph{almost Dold}.
An almost Dold sequence with~$\least_n(a)\ge0$
for all~$n\in\NN$ is called \emph{almost realizable}.
\end{definition}

This property leads to the notion of `repairing' a sequence
in the following sense. If~$\failure(a)$ exists (that is, is finite)
then multiplying~$a$ by~$\failure(a)$ produces a Dold sequence
and~$\failure(a)$ is the smallest integer with this property.
We will see later that the Lucas sequence gives rise
to an infinite family of non-trivial almost Dold sequences
to sit alongside the Stirling numbers identified in~\cite{MR4361581}.

\begin{definition}[Good primes]
Given an integer sequence~$a$ define two
associated sets of
primes as follows:
\begin{align*}
\gooddoldprimes(a)
&=
\{p\in\PP :
a\mbox{ is Dold at }p\},\\
\goodrealizableprimes(a)
&=
\{p\in\PP :
a\mbox{ is realizable at }p\}
\end{align*}
and similarly define~$\gooddoldprimes^{\text{almost}}$ and~$\goodrealizableprimes^{\text{almost}}$
for the analogous sets using the almost Dold and almost realizable
properties.
\end{definition}

We will also write~$\baddoldprimes(a)=\PP\setminus\gooddoldprimes(a)$
and~$\badrealizableprimes(a)=\PP\setminus\goodrealizableprimes(a)$
for the bad primes.
In the paper~\cite{MR5076929} these sets were discussed
for the Bernoulli denominators, the Bernoulli
numerators, and the Euler numbers in the particular
context of understanding realizability by a group
automorphism rather than an arbitrary map.
In those settings~$\gooddoldprimes(a)$ can be
viewed as a generalization of the notion of
regular prime (for the Bernoulli denominators)
for other sequences.

Clearly there are degenerate sequences for which all
primes are \emph{good},~$(1,1,1,\dots)$ being an example,
but the same phenomenon can arise in non-trivial ways.

\begin{example}[All primes can be good]
In~{\rm{\cite[Th.~4]{MR5076929}}} it is shown that an automorphism of
a nipotent group is locally nilpotently realizable at every
prime. That is, if~$\theta\colon G\to G$ is an automorphism
of a nilpotent group with the property that the
sequence~$a=(\fix_n(\theta))_{n\in\NN}$ defined by
\[
\fix_n(\theta)
=
\vert\{g\in G\colon\theta^ng=g\}\vert
\]
is finite for all~$n\in\NN$, then for any~$p\in\PP$
there is an automorphism~$\theta_p$ of a locally
nilpotent group~$G_p$ for which
\[
\fix_n(\theta_p)=[\fix_n(\theta)]_p
\]
for all~$n\in\NN$. That is,~$\goodrealizableprimes(a)=\PP$.
\end{example}

It is equally clear that there are degenerate sequences
for which all primes are \emph{bad}.
If~$a=(n)_{n\in\NN}$, then for any~$p\in\PP$ we have~$a_p=p$
and so~$a_p-a_1$ is not congruent to~$0$ modulo~$p$.
It follows that~$\baddoldprimes(a)=\PP$.

Our purpose here is to discuss these sets of
primes and their `almost' analogues for certain linear recurrence sequences.
We deliberately isolate some of the arguments for the Dold
property on its own, as this is the natural property for
arithmetic. Realizability is natural for dynamical systems,
but involves sign questions that are not natural in
arithmetic.
The picture that emerges is far from complete---indeed we
have nothing to say beyond a certain family of
binary recurrences---but is we hope of some interest.

\section{Lucas and Fibonacci numbers}

We will consider the Lucas sequence~$L=(1,3,4,\dots)$
and its companion the Fibonacci sequence~$F=(1,1,2,\dots)$
throughout this section. Define the two ranks of apparition
of a prime~$p$ by
\begin{align*}
\lambda(p)
&=
\begin{cases}
\min\{n\in\NN : p\divides L_n\}&\mbox{if $p$ divides some term of $L$},\\
\infty&\mbox{if not}
\end{cases}
\intertext{and}
\varphi(p)
&=
\begin{cases}
\min\{n\in\NN : p\divides F_n\}&\mbox{if $p$ divides some term of $F$},\\
\infty&\mbox{if not.}
\end{cases}
\end{align*}
It is well-known that
\begin{equation*}\label{rank1ForFibonacci}
\varphi(p)
\divides
p-\Bigl(\frac{5}{p}\Bigr)
\end{equation*}
for any odd prime~$p\neq5$, where
we write~$\bigl(\frac{5}{p}\bigr)$ for the Legendre symbol.

Clearly if~$p$ is a prime that does not divide
any term of~$L$ then~$p\in\goodrealizableprimes(L)$, so
we first examine those. Throughout we will make use of
well-known congruences and rules for rank of apparition
for the Lucas and Fibonacci sequences as they may be found, for example,
in Ribenboim's survey~\cite{MR1352481}.

\begin{lemma}\label{lemmaTrivialLucas}
Let~$p$ be a prime. Then the values of~$\lambda$ are
as follows.
\begin{enumerate}
\item[\rm(1)] $\lambda(2)=3$.
\item[\rm(2)] $\lambda(5)=\infty$.
\item[\rm(3)] If~$p\ne5$, then~$\lambda(p)<\infty$ if and only
if~$\varphi(p)$ is even. More precisely,
\begin{equation*}\label{eq:twolambda}
 \varphi(p)=2\lambda(p).
 \end{equation*}
\item[\rm(4)] If~$p\ne5$ and~$\left(\frac5p\right)=-1$, then
we have
\begin{equation}\label{oddrank}
\varphi(p)\text{ is odd}\quad\Longleftrightarrow\quad p\equiv1\pmod4.
\end{equation}
Consequently, if~$p\neq5$ and~$\left(\frac5p\right)=-1$ then
\[
\lambda(p)=\infty
\Longleftrightarrow
p\equiv13,17\pmod{20}
\]
and
\[
\lambda(p)
<
\infty
\Longrightarrow
p\equiv3,7\pmod{20}.
\]
\end{enumerate}
\end{lemma}

\begin{proof}
Firstly notice that~$2\divides L_3$,
proving~(1).
Now~$L_n=\frac{F_{2n}}{F_{n}}$
and~$[F_{2n}]_5=[F_n]_{5}$ so~$5$
does not divide any term of~$L$,
giving~(2).
If~$p\neq2,5$ then~$p\divides L_n$ if and only
if~$p\divides F_{2n}$ and~$p\notdivides F_{n}$.
It follows that~$p$ divides some term of~$L$
if and only if~$\varphi(p)$ is even.
Since~$p\divides F_m$ if and only if~$\varphi(p)\divides m$,
this becomes $\varphi(p)\divides 2n$ and $\varphi (p)\notdivides n$.
Such an index~$n$ exists precisely when~$\varphi(p)$ is even,
and the least such~$n$ is~$\frac{\varphi(p)}{2}$,
which shows~(3).
For~(4), assume that~$p\ne2,5$ and~$\left(\frac5p\right)=-1$
so that~$p\equiv\pm 2$ modulo~$5$
and~$5$ is a quadratic non-residue modulo~$p$ and
hence~$\FF_p[\sqrt{5}]=\FF_{p^2}$.
Let~$\alpha=\frac{(1+\sqrt5}{2}$,~$\beta=\frac{(1-\sqrt5)}{2}$,
and~$x=\frac{\alpha}{\beta}\in\mathbb F_{p^2}$.
The Frobenius exchanges~$\alpha$ and $\beta$,
so $x^p=x^{-1}$,~$\varphi(p)\divides p+1$,
and~$x=-\alpha^2$ so
\[
 x^{(p+1)/2}=(-1)^{(p+1)/2}\alpha^{p+1}=(-1)^{(p+3)/2}.
\]
If~$p\equiv1\pmod4$, then the right-hand side is~$1$
and $\frac{p+1}{2}$ is odd, so~$\varphi(p)$ is odd.
If~$p\equiv3\pmod4$, then the value~is~$-1$, so the order of~$x$ is even
and~$\varphi(p)$ is even and~\eqref{oddrank}
follows.
The residue-class consequences now follow by the Chinese remainder theorem.
\end{proof}

\begin{lemma}\label{lemmaMainLucas}
For~$p\in\PP$ with~$p\neq2,3$ we have~$p\in\gooddoldprimes(L)$ if
and only if~$\lambda(p)=\infty$ or~$2\lambda(p)\divides p-1$.
The primes~$2$ and~$3$ are not in~$\gooddoldprimes(L)$.
\end{lemma}

\begin{proof}
Write~$a_n=[L_n]_p$ for all~$n\ge1$ and~$p\in\mathbb{P}$
and~$r=\lambda(p)$ for~$p\in\mathbb{P}$ throughout,
with the prime~$p$ being clear from context as we go through.

For~$p=2$ a direct calculation gives~$
\least_6([L]_2)=a_1-a_2-a_3+a_6=-2$,
which is not divisible by~$6$.
Similarly for~$p=3$ we have~$\least_4([L]_3)=-a_2+a_4=-3+1=-2$.
which is not divisible by~$4$, so~$2,3\notin\gooddoldprimes(L)$.

Now assume that~$p\in\mathbb{P}\setminus\{2,3\}$.
We know that~$p\divides L_n$ if and only if~$n\in\lambda(p)(2\NN_0+1)$,
indeed we have the following `lifting the exponent' result for
the Lucas sequence:
\begin{equation*}\label{equationPAdicSizeLucasSequences}
\ord_p(L_{\lambda(p)k})
=
\ord_p(L_{\lambda(p)})+\ord_p(k)
\end{equation*}
for~$k\in2\NN_0+1$.
Thus by induction we have
\begin{equation}\label{equationPAdicSizeLucasSequences2}
[L_n]_p
=
\begin{cases}
p^{\ord_p(L_{\lambda(p)}+\ord_p(k))}&\mbox{if }n\in\lambda(p)k\mbox{ with $k$ odd};\\
1&\mbox{if not.}
\end{cases}
\end{equation}
Suppose now that~$p\in\gooddoldprimes(L)$ and
let~$c=\ord_p(a_r)$ so~$c>0$ and~$a_d=1$
for~$d<r$, for~$d\in2r\mathbb{N}$, and for~$d$ not a multiple of~$r$.

We first claim that every odd prime-power factor of~$2r$
divides~$p-1$.
To see this, write~$r=q^st$ with~$q\in\mathbb{P}\setminus\{2\}$
and~$q\notdivides t$ and let~$m=tp^j$ for some~$j\ge0$.
Then~$q\notdivides m$ since~$p\notdivides r$,
so~$mq^s=rp^j$ is an odd multiple of~$r$,
and hence~$a_{mq^s}=p^{c+j}$ by~\eqref{equationPAdicSizeLucasSequences2}.
By~\eqref{equationDoldForIntegerSequence2} we deduce
that~$p^{c_j}\equiv1$ modulo~$q^s$ and therefore in particular
\[
p^{c+1}-p^c=p^c(p-1)\equiv0\pmod{q^s}
\]
which shows that
\begin{equation}\label{equationcleveroldSawian}
q^s\divides p-1
\end{equation}
as claimed.

Now assume that~$r=2^st$ with~$s>0$ and~$t$ odd,
and write~$m=tp^j$ with~$j\ge0$. Since~$p\neq2$
we know that~$m$ is odd and so~$m2^s=rp^j$ is
once again an odd multiple of~$r$.
It follows that~$a_{m2^s}=p^{c+j}$.
On the other hand~$m2^{s+1}$ is an even multiple
of~$r$ and so~$a_{m2^{s+1}}=1$. By~\eqref{equationDoldForIntegerSequence2}
we deduce that~$p^{c+j}\equiv1$ modulo~$2^{s+1}$
and so in particular
\[
p^{c+1}-p^c=p^c(p-1)\equiv0\pmod{2^{s+1}}.
\]
Thus~$2^{s+1}\divides p-1$ which, combined
with~\eqref{equationcleveroldSawian} for every odd~$q$,
shows that~$2r\divides p-1$.

We now turn to the reverse direction.
If~$r$ does not exist (that is, if~$p$ does not divide
any term of~$L$) then~$p\in\gooddoldprimes(L)$
trivially. So suppose that~$r$ exists and~$2r\divides p-1$
and~$q\notdivides m$.
If~$m$ is not an odd multiple of~$r$ then
\[
a_{mp^s}=a_{mp^{s-1}}=1
\]
by~\eqref{equationPAdicSizeLucasSequences2}.
If~$m=rk$ for some odd~$k$ then~$p\notdivides k$ since~$p\notdivides m$,
and so
\[
a_{mp^s}-a_{mp^{s-1}}=p^{c+s}-p^{c+s-1}=p^{c+s-1}(p-1)\equiv0\pmod{p^s},
\]
which is compatible
with~\eqref{equationDoldForIntegerSequence2}.
If~$q\neq p$ and~$r\divides m$ then~$r\divides mq^s$ for any~$s\ge1$
and so
\[
\ord_{p}(mq^s/r)=\ord_p(m/r).
\]
It follows that~$a_{mq^s}\equiv q_{mq^{s-1}}$ modulo~$q^s$,
which again is compatible with~\eqref{equationDoldForIntegerSequence2}.
If~$r\notdivides m$ then the only way in which multiplication
by~$q$ can make~$r$ divide~$mq^s$ is for~$q$ to complete
the~$q$-part of~$r$. However,~$2r\divides p-1$ so~$p\equiv1$
modulo~$q^s$ for those~$s$ where this is possible.
It follows that~$a_{mq^s}\equiv a_{mq^{s-1}}$ modulo~$q^s$,
which completes the proof of~\eqref{equationDoldForIntegerSequence2}
and hence shows that~$p\in\gooddoldprimes(L)$.
\end{proof}



\begin{theorem}\label{theoremLucas}
The Lucas sequence~$L$ is almost Dold at every
prime, and
\[
\failure([L]_p)
=
\begin{cases}
3&\mbox{if }p=2,\\
\lambda(p)&\mbox{if }q\equiv3,7\pmod{20},\\
1&\mbox{in all other cases}.
\end{cases}
\]
\end{theorem}

\begin{proof}
Assume that~$p\neq2,3$.
We first prove that~$L$ is Dold at a prime
if and only if~$p\neq2$ and~$p$ is not congruent to~$3$
or~$7$ modulo~$20$, giving the third case.

If~$p\equiv\pm1$ modulo~$5$ then~$2\lambda(p)\divides p-1$,
and we may apply Lemma~\ref{lemmaMainLucas}
to see that~$p\in\gooddoldprimes(L)$.

If~$p\equiv\pm2$ modulo~$5$ then~$\varphi(p)\divides p+1$
and there are two cases to consider.
If~$p\equiv3$ modulo~$4$ then the Fibonacci rank~$\varphi(p)$
is even and hence~$p$ divides a Lucas number
and~$2\lambda(p)=\varphi(p)$.
The Dold condition requires that~$2\lambda(p)=\varphi(p)\divides p-1$
which is impossible as~$\varphi\divides p+1$ and~$\varphi(p)>2$
by inspection. Thus the simultaneous congruences~$p\equiv3$
modulo~$4$ and~$p\equiv\pm2$ modulo~$5$ imply that~$p\in\baddoldprimes(L)$.
By the Chinese remainder theorem these bad primes
are precisely those congruent to~$3$ or to~$7$ modulo~$20$.
If~$p\equiv1$ modulo~$4$ and~$p\equiv\pm2$ modulo~$5$
then Lemma~\ref{lemmaTrivialLucas} shows that~$p
\in\gooddoldprimes(L)$ trivially.

Applying the Chinese remainder theorem and checking
off the possible congruence classes modulo~$20$ then
proves that~$\failure([L]_p)=1$ unless~$p=2$ or~$p\equiv3,7$ modulo~$20$
and proves that~$\failure([L]_p)$ is not~$1$ (and
{\it{a priori}} may not exist) in all the other
cases.

The calculation of the claimed repair factors
for~$p=2$ and for~$p\equiv3,7$ modulo~$20$
involves precise knowledge of all the values
of~$\least_n([L]_p)$.
By Lengyel~\cite[Lemma~2]{MR1337793}
\[
[L_n]_2
=
\begin{cases}
1&\mbox{if }n\equiv1,2\pmod{3},\\
2&\mbox{if }n\equiv0\pmod{6},\\
4&\mbox{if }n\equiv3\pmod{6}
\end{cases}
\]
and so
\[
\least_n([L]_2)
=
\begin{cases}
1&\mbox{if }n=1,\\
3&\mbox{if }n=3,\\
-2&\mbox{if }n=6,\\
0&\mbox{for all other $n$}.
\end{cases}
\]
Thus~$\failure([L]_2)=3$ as claimed.

A direct calculation
using~\eqref{equationPAdicSizeLucasSequences2} shows
that
\begin{equation}\label{equationAA}
\least_{n}([L]_p)
=
\begin{cases}
\displaystyle\sum_{\substack{d\smalldivides m\\d\text{ odd}}}
\mu(\tfrac{m}{d})(p^{e+\ord_p(d)}-1)
&\mbox{if }n=rm,\\
0&\mbox{if }r\notdivides n.
\end{cases}
\end{equation}

If~$\lambda(p)$ is odd then~$[L_n]_p=1$ for all~$n\ge1$
and we are in the good prime case.
So assume that~$p\equiv3,7$ modulo~$20$
and~$\lambda(p)=r$.
We claim that if~$e=\ord_p(L_{r})$ the
only non-zero terms of~$\least_n([L]_p)$ are
\begin{align}
\least_1([L]_p)&=1,\nonumber\\
\least_{r}([L]_p)&=p^e-1\mbox{ if }r>1,\label{was12}\\
\least_{2r}([L]_p)&=1-p^e,\nonumber\\
\least_{rp^j}([L]_p)&=p^{e+j-1}(p-1)\mbox{ for }j\ge1,\nonumber\\
\least_{2rp^j}([L]_p)&=-p^{e+j-1}(p-1)\mbox{ for }j\ge1\label{was15}
\end{align}
by~\eqref{equationAA}.
This gives the remaining statement of the theorem
since it shows that~$p$
lies in~$\gooddoldprimes(L)$
if and only if~$\lambda(p)\divides p-1$ and that
\[
\failure([L]_p)
=
\frac{\lambda(p)}{\gcd(\lambda(p),p-1)}.
\]

If~$m=2^as$ with~$s$ odd then~\eqref{equationAA}
vanishes if~$a\ge2$ and for~$a=1$ is minus
its value at~$s$.
If~$a=0$ and~$r>1$ then~\eqref{equationAA} is~$p^e-1$
when~$s=1$ and for~$s>1$ we have
\begin{equation}\label{equationAB}
\least_{rs}([L]_p)
=
\sum_{d\smalldivides s}\mu(\tfrac{s}{d})a_d.
\end{equation}
The claimed equations~\eqref{was12}--\eqref{was15}
now follow from~\eqref{equationAA} and~\eqref{equationAB}.
\end{proof}

\begin{thebibliography}{99}
\bibitem[MR1337793]{MR1337793} T.~Lengyel. \newblock The order of the {F}ibonacci and {L}ucas numbers. \newblock {\em Fibonacci Quart.}, 33(3):234--239, 1995.
\bibitem[MR1352481]{MR1352481} P.~Ribenboim. \newblock The {F}ibonacci numbers and the {A}rctic {O}cean. \newblock In {\em Proceedings of the 2nd {G}auss {S}ymposium. {C}onference {A}: {M}athematics and {T}heoretical {P}hysics ({M}unich, 1993)}, Sympos. Gaussiana, pages 41--83. de Gruyter, Berlin, 1995.
\bibitem[MR3195758]{MR3195758} G. T. Minton. \newblock {Linear recurrence sequences satisfying congruence conditions}. \newblock {\em Proc. Amer. Math. Soc.}, {142}(7): 2337--2352, 2014,
\bibitem[MR4332826]{MR4332826} J.~Byszewski, G.~Graff, and T.~Ward. \newblock Dold sequences, periodic points, and dynamics. \newblock {\em Bull. Lond. Math. Soc.}, 53(5):1263--1298, 2021.
\bibitem[MR4361581]{MR4361581} P.~Miska and T.~Ward. \newblock Stirling numbers and periodic points. \newblock {\em Acta Arith.}, 201(4):421--435, 2021.
\bibitem[MR5076929]{MR5076929} S.~Jaidee, P.~Moss, and T.~Ward. \newblock Local structure of classical sequences, regular primes, and dynamics. \newblock {\em Res. Number Theory}, 12(2):Paper No. 54, 27, 2026.
\end{thebibliography}
\end{document}
